\documentclass[10pt]{article}
\usepackage[a4paper,margin=2.1cm]{geometry}
\usepackage{amsmath,amssymb,booktabs,graphicx,enumitem}
\usepackage[T1]{fontenc}\usepackage{lmodern}
\usepackage[font=small]{caption}
\setlength{\parindent}{0pt}\setlength{\parskip}{5pt}
\newcommand{\bs}{\beta^\star}\newcommand{\E}{\mathbb E}\newcommand{\asinh}{\operatorname{arcsinh}}\newcommand{\prox}{\operatorname{prox}}
\title{One template and its variants:\\ cross-entropy, early stopping, SGD and regularisers in the diagonal-network model}
\author{}\date{23 November 2026}
\begin{document}\maketitle

\section{The template}
The model is $f(x)=x^\top\beta$ with $\beta=u^2-v^2$. Data: $n$ inputs $x_i\sim\mathcal N(0,I_d/d)$, the columns of $X$, with targets $y$; $\alpha=n/d$. Teacher $\bs$: entries $\pm\rho^{-1/2}$ with probability $\rho$, else $0$.

\textbf{(T1) Mirror flow.} For any loss $L(X^\top\beta)$, gradient flow on $(u,v)$ is
\begin{equation}\label{T1}
\beta=c\,\sinh\theta,\qquad \dot\theta=-4\,\nabla_\beta L ,\qquad c=2uv\ \text{(constant in time)}.
\end{equation}
So $\theta(t)-\theta(0)$ always lies in the span of the inputs. Small $c$ is the rich regime, large $c$ the lazy one ($c=\sqrt k/2$).

\textbf{(T2) End point.} If training from $\beta_0$ ends with $X^\top\beta=y$, then
\begin{equation}\label{T2}
\beta_\infty=\arg\min_\beta D_c(\beta,\beta_0)\quad\text{s.t.}\quad X^\top\beta=y ,
\end{equation}
where $D_c$ is the Bregman divergence of $\psi_c$, $\psi_c'(\beta)=\asinh(\beta/c)$. (These are the KKT conditions of \eqref{T2}: $\theta_\infty-\theta_0\in\operatorname{span}X$ and the constraints.)

\textbf{(T3) Large $d$.} For a problem ``minimise $\sum_jR(\beta_j)$ subject to $X^\top\beta=X^\top\bs$'', each weight is distributed as
\begin{equation}\label{T3}
\beta=\prox_{\tau R}(\bs+\sigma h),\quad h\sim\mathcal N(0,1),\qquad \sigma^2=\frac{\E(\beta-\bs)^2}{\alpha},\qquad \E[\prox{}']=\alpha .
\end{equation}
Two unknown numbers, $\tau$ and $\sigma$. Pretraining from zero is $R=\psi_c$.

Every variant below changes one of three things: the \emph{scale} $c$, the \emph{start} $\beta_0$, or the \emph{constraints}. Table~\ref{tab} is the summary.

\begin{table}[h]\centering\small
\begin{tabular}{@{}lll@{}}\toprule
variant & what changes in the template & status\\\midrule
cross-entropy, soft labels & targets $y\to$ label logits & exact\\
rescaling the weights by $\sqrt\kappa$ & $c\to\kappa c$, $\beta_0\to\kappa\beta_0$ & exact\\
weight decay, then switched off & the same, with $\kappa=e^{-2\Lambda}$ & exact\\
mini-batch SGD & the same, with $\kappa=\prod_k(1-4\gamma^2g_k^2)$ & exact; $\kappa$ random\\
one-sided decay, then switched off & $c\to e^{-\Lambda}c$ and $\theta_0\to\theta_0-\Lambda$ & exact\\
regulariser kept on & \eqref{T2} becomes a penalised fit; start forgotten & exact (convex $R$)\\
early stopping & targets $\to$ current predictions; needs dynamics & exact, solved numerically\\
ascent on hard labels & no end point & approximate only\\
non-convex regulariser & end point depends on the path & no formula\\\bottomrule
\end{tabular}\caption{The variants.}\label{tab}\end{table}

\section{Cross-entropy}
\textbf{Soft labels.} If the labels are probabilities $p_i\in(0,1)$, the loss is minimised at the finite logits $\log\frac{p_i}{1-p_i}$. By \eqref{T1} the flow converges to \eqref{T2} with those logits as targets. With a logistic teacher, $p_i=\sigma(x_i^\top\bs)$, the targets are $x_i^\top\bs$: \emph{cross-entropy pretraining is the squared-loss model}, and everything known for it carries over. (Flow vs \eqref{T2}: $2\times10^{-6}$.)

\textbf{Hard labels.} The weights diverge and only a limiting direction is known (Moroshko et al.; Pesme et al.). At realistic losses neither the hard-margin projection ($27$--$54\%$ off) nor the regularisation path ($19$--$38\%$) describes the flow. We do not use hard labels for pretraining.

\textbf{Ascent on hard labels} (gradient ascent, NPO). With $z_i=y_ix_i$ and per-sample weight $w(\cdot)>0$ on the margin $m_i=z_i^\top\beta$,
\begin{equation}
\theta(t)=\theta_0-\sum_i\Lambda_i(t)\,z_i,\qquad \dot\Lambda_i=4\,w(m_i).
\end{equation}
The losses differ only in $w$ (ascent: $\sigma(-m)$; NPO: that times a factor that decays as the sample is forgotten). There is no end point. Freezing the $\Lambda_i$ at a common value is accurate to $3$--$8\%$ early and $14$--$35\%$ near chance. In simulation, for every $k$ from $10^{-8}$ to $10^{4}$, ascent gains at most $6\%$ over doing nothing, while soft targets at the uniform label (zero logit, exact by \eqref{T2}) do far better.

\section{Rescaling, weight decay, SGD: one parameter $\kappa$}
\textbf{Rescaling.} Multiplying all weights by $\sqrt\kappa$ gives $c\to\kappa c$, $\beta_0\to\kappa\beta_0$, $\theta_0$ unchanged. Then \eqref{T2} applies with the new scale and start.

\textbf{Weight decay that is switched off.} With penalty $\frac{\lambda_t}2(\|u\|^2+\|v\|^2)$,
$\dot u=-2u\,\nabla L-\lambda_tu$, $\dot v=2v\,\nabla L-\lambda_tv$, hence
\begin{equation}
\beta_t=c\,e^{-2\Lambda_t}\sinh\theta_t,\qquad\dot\theta=-4\nabla_\beta L,\qquad\Lambda_t=\int_0^t\lambda_s\,ds .
\end{equation}
The decay never touches $\theta$; it only shrinks the scale. After switch-off and convergence the end point is \eqref{T2} with $\kappa=e^{-2\Lambda}$: \emph{temporary weight decay is a rescaling}, and only its total $\Lambda$ matters, not its timing. This is Theorem 3.11 of Jacobs, Zhou \& Burkholz (2025) in our model; their ``type of bias'' is $c\to\kappa c$ and their ``positional bias'' is $\beta_0\to\kappa\beta_0$.

\textbf{One-sided decay} $\frac{\lambda_t}2\|u\|^2$ (their Example 3.13): $u$ shrinks by $e^{-\Lambda}$ and $v$ does not, so $c\to e^{-\Lambda}c$ and $\theta_0\to\theta_0-\Lambda$. For pretraining this is \eqref{T3} with $R(\beta)=\psi_{e^{-\Lambda}c}(\beta)+\Lambda\beta$.

\textbf{Mini-batch SGD.} One step with step size $\gamma$ and mini-batch gradient $g_k$ (in $\beta$) multiplies $u$ by $1-2\gamma g_k$ and $v$ by $1+2\gamma g_k$. So $uv$ shrinks by $1-4\gamma^2g_k^2$ per step, and to $O(\gamma^3)$ the mirror variable moves as in \eqref{T1}:
\begin{equation}
\kappa_j=\prod_k\big(1-4\gamma^2g_{k,j}^2\big)\qquad\text{(Even et al.\ 2023).}
\end{equation}
\emph{SGD is gradient flow with a smaller scale}; in stage 2 it is a rescaling of the pretrained model. $\kappa$ depends on the trajectory. For batches of size $b$ out of $n$, per step $\E\,g_j^2=\bar g_j^2+c_bV_j$, with $\bar g$ the full gradient, $V_j$ the per-sample variance and $c_b=\frac{n-b}{b(n-1)}$. In expectation, to first order in $\gamma$,
\begin{equation}
\E\log\frac1{\kappa_j}=4\gamma\int_0^\infty\!\big(c_bV_j+\bar g_j^2\big)\,dt ,
\end{equation}
the integrals taken along gradient flow. The relative error is of order $\gamma c_b$ ($11\%$ at $\gamma c_b=0.02$, $2\%$ at $0.005$); it comes from the path changing with the noise and has no closed form. As $\gamma\to0$, $\kappa\to1$ for every batch size: SGD with infinitesimal steps is gradient flow.

\section{A regulariser that stays on}
With a penalty $\lambda\sum_jR(\beta_j)$ kept on, training converges to the minimiser of loss plus penalty. For convex $R$ it is unique: \emph{the start and the scale are forgotten}. Weight decay on $(u,v)$ is $R=|\beta|$ (the cheapest factorisation of $\beta$ costs $|\beta|$), so the end point is a lasso. Large $d$: \eqref{T3} with one change,
\begin{equation}\label{reg}
\beta=\prox_{\tau R}(\bs+\sigma h),\qquad \sigma^2=\frac{\E(\beta-\bs)^2}{\alpha},\qquad \E[\prox{}']=\alpha-\frac\lambda\tau .
\end{equation}
($\lambda\to0$ recovers \eqref{T3}.) With several groups of samples with different targets, $\bs$ and $\alpha$ are replaced by group-weighted sums; the formulas are in \texttt{se\_wd.py}.

For non-convex $R$ (weight decay on a three-layer diagonal network gives $|\beta|^{2/3}$) this fails: different algorithms and different initialisations end at different points (norm between $0.14$ and $0.30$ in one setting), and the solution of \eqref{reg} ($0.42$) is none of them. The end point is path-dependent.

\section{Early stopping}
\textbf{Identity.} By \eqref{T1}, at every time $\beta(t)$ is \eqref{T2} with targets $X^\top\beta(t)$: early stopping is the same bias applied to partly fitted targets. What is missing is how the predictions move.

\textbf{Decoupled inputs.} If each coordinate has its own target $b$ and input scale $s$, then $\dot\beta=-4s\sqrt{\beta^2+c^2}\,(\beta-b)$, solved by
\begin{equation}
\frac{e^{\theta}-e^{\theta_b}}{e^{\theta}+e^{-\theta_b}}=\text{const}\cdot e^{-4s\sqrt{b^2+c^2}\,t},\qquad \sinh\theta_b=b/c .
\end{equation}
Forgetting to $b=0$ at small $c$ is a power law, $\beta(t)\approx\beta_0/(1+4s\beta_0t)$, and what is left relearns faster than from scratch ($0.37$ against $1.12$ from a leftover of $0.1$): stopped early, a weight is hidden, not deleted. At large $c$ both are exponential and equal.

\textbf{Gaussian inputs, large $d$.} Let the targets be $a(t)\,y$ (so pretraining, then a switch of $a$, is one flow). Each coordinate obeys
\begin{equation}
\dot\theta=-4\Big(\eta+\alpha\,K(\beta-a\bs)\Big),\qquad \beta=c\sinh\theta,
\end{equation}
where $\eta$ is a Gaussian process and $K=(1-R)^{-1}$ a causal operator in time, with $R(t,s)=\E\,\partial\beta(t)/\partial\eta(s)$. The residual of a training sample is $K\zeta$, $\zeta$ Gaussian with the covariance of $\beta-a\bs$, and $\eta=\sqrt\alpha\,K\zeta$. Groups of samples with different targets or weights add their terms. This extends Nishiyama \& Imaizumi (2025) from a zero start to a switch. It is solved by one forward pass in time over sampled paths. A second copy trained on the kept samples only, driven by the same randomness, gives the distance to the retrained model at every stopping time.

Its sampling noise is as large as a simulation's of the same size: at $k=10^{-6}$ single runs with 1000--2000 paths scatter by $\pm0.02$ to $\pm0.05$ on an untraining distance of $0.19$, which is what looked like a $10\%$ gap earlier. Averaged over runs it agrees with simulation ($0.177\pm0.013$ against $0.188$).

\section{Checks}
\begin{center}\small\begin{tabular}{@{}p{6.1cm}p{9.2cm}@{}}\toprule
claim & check\\\midrule
soft-label cross-entropy $=$ \eqref{T2} & flow vs projection, $2\times10^{-6}$, three scales\\
its large-$d$ curves against $k$ & 14 settings, within one s.e.\ ($d=400$, 4 seeds)\\
switched-off weight decay $=$ rescaling & gradient descent vs \eqref{T2}: $0.1$--$1.2\%$ (step-size error); same end point for decay early or late\\
its large-$d$ curves & 42 numbers, $d=1600$, 16 seeds: all within 2 s.e.\\
one-sided decay, large $d$ & 36 numbers, $d=1600$, 24 seeds: all within 1 s.e.\\
regulariser kept on, three choices of $R$ & 54 numbers, $d=1500$ and $6000$: all within 2.2 s.e.\\
SGD $=$ flow from the shrunk scale & $10^{-4}$ to $4\times10^{-3}$; stage 2: $0.2\%$ against $4.7\%$ without $\kappa$\\
SGD batch-size formula & mean of 8--16 runs: ratio $1.11,1.05,1.02$ at $\gamma=0.02,0.01,0.005$ ($b=1$)\\
early-stopping dynamics & norm and overlap through both stages at $k=1,10^{-2},10^{-6}$; reproduces the optimal stopping time at small $k$\\\bottomrule
\end{tabular}\end{center}
Across the three regulariser checks: 132 numbers, $98\%$ within 2 s.e., largest $2.2$ (Figure~\ref{fig}).

\begin{figure}[h]\centering\includegraphics[width=\linewidth]{fig_regularisers_tight.png}
\caption{Regulariser checks. Lines: theory. Markers: exact finite-size end points, $\pm1$ s.e.}\label{fig}\end{figure}

\section*{Open}
\begin{itemize}[nosep,leftmargin=1.2em]
\item Ascent losses (gradient ascent, NPO): no formula; the dynamical theory with a nonlinear sample side would give one.
\item SGD beyond first order in $\gamma$, and multi-pass mini-batches at large $d$.
\item Non-convex regularisers; deeper networks.
\item Stage 2 with one-sided decay, and distances between two regularised solutions on overlapping data.
\end{itemize}

\small\textbf{References.} Woodworth et al., COLT 2020. Moroshko et al., NeurIPS 2020. Pesme, Pillaud-Vivien, Flammarion, NeurIPS 2021. Even, Pesme, Gunasekar, Flammarion, NeurIPS 2023. Pesme, Dragomir, Flammarion, NeurIPS 2024. Jacobs, Zhou, Burkholz, ICML 2025 (arXiv:2504.12883). Nishiyama, Imaizumi, arXiv:2510.01930.
\end{document}
